% !TeX program = tectonic
\documentclass[11pt]{article}
\IfFileExists{alius-guidelines-style.tex}{\usepackage[a4paper,top=27mm,bottom=24mm,left=25mm,right=25mm]{geometry}
\usepackage[T1]{fontenc}
\usepackage[utf8]{inputenc}
\usepackage[scaled=0.96]{helvet}
\renewcommand{\familydefault}{\sfdefault}
\usepackage{array}
\usepackage{enumitem}
\usepackage{fancyhdr}
\usepackage[hidelinks]{hyperref}
\usepackage{microtype}
\usepackage{needspace}
\usepackage{tabularx}
\usepackage{xcolor}
\usepackage{xurl}

\definecolor{aliusgreen}{HTML}{13823A}
\definecolor{aliusgray}{HTML}{7A7A7A}
\definecolor{aliusline}{HTML}{9A9A9A}

\hypersetup{
  colorlinks=true,
  linkcolor=aliusgreen,
  urlcolor=aliusgreen
}

\setlength{\parindent}{0pt}
\setlength{\parskip}{0.72em}
\setlength{\tabcolsep}{0pt}
\raggedbottom
\sloppy
\urlstyle{same}

\newcommand{\ALIUSFooterTitle}{ALIUS Bulletin Guidelines (2022)}

\pagestyle{fancy}
\fancyhf{}
\fancyfoot[L]{\footnotesize\color{aliusgray}\ALIUSFooterTitle}
\fancyfoot[R]{\footnotesize\color{aliusgray}\thepage}
\renewcommand{\headrulewidth}{0pt}
\renewcommand{\footrulewidth}{0.45pt}

\setlist[itemize]{
  leftmargin=1.35em,
  itemsep=0.18em,
  topsep=0.15em,
  parsep=0em
}
\setlist[enumerate]{
  leftmargin=2.05em,
  itemsep=0.42em,
  topsep=0.2em,
  parsep=0em
}
\setlist[description]{
  leftmargin=0em,
  labelsep=0.55em,
  itemsep=0.35em,
  topsep=0.2em,
  font=\normalfont\bfseries
}

\newcommand{\guideDivider}{%
  \par\noindent{\color{aliusline}\rule{\textwidth}{0.45pt}}\par
}

\newcommand{\guideMeta}[2]{%
  {\footnotesize\bfseries\MakeUppercase{#1}}\par
  {\footnotesize #2}\par\vspace{0.65em}
}

\newcommand{\guideHeader}[5]{%
  \thispagestyle{fancy}%
  \vspace*{1mm}
  \begin{center}
    {\Large\bfseries\MakeUppercase{ALIUS Bulletin}\par}
    \vspace{0.25em}
    {\normalsize *\par}
    \vspace{0.25em}
    {\normalsize\bfseries\MakeUppercase{Guidelines For}\par}
    \vspace{0.15em}
    {\LARGE\bfseries\MakeUppercase{#1}\par}
  \end{center}
  \vspace{1.05em}
  \begin{tabularx}{\textwidth}{@{}>{\raggedright\arraybackslash}p{0.55\textwidth}!{\hspace{1.35em}{\color{aliusline}\vrule width 0.5pt}\hspace{1.35em}}>{\raggedright\arraybackslash}X@{}}
    {\large\itshape #2} &
    \guideMeta{Format}{#3}
    \guideMeta{Editorial Route}{#4}
    \guideMeta{Contact}{#5}
  \end{tabularx}
  \vspace{0.75em}
  \guideDivider
  \vspace{0.25em}
}

\newcommand{\guideSection}[1]{%
  \Needspace{6\baselineskip}%
  \par\vspace{1.1em}
  \begin{center}
    {\color{aliusline}\rule{0.14\textwidth}{0.45pt}}\hspace{0.85em}%
    {\large\bfseries #1}%
    \hspace{0.85em}{\color{aliusline}\rule{0.14\textwidth}{0.45pt}}
  \end{center}
  \vspace{-0.25em}
}

\newcommand{\guideSubsection}[1]{%
  \par\vspace{0.65em}{\bfseries #1}\par\vspace{-0.35em}
}

\newcommand{\guideQuestion}[1]{%
  \par\vspace{0.55em}{\color{aliusgreen}\bfseries #1}\par
}

\newenvironment{guideChecklist}
  {\begin{itemize}[label=\textendash]}
  {\end{itemize}}
}{\usepackage[a4paper,top=27mm,bottom=24mm,left=25mm,right=25mm]{geometry}
\usepackage[T1]{fontenc}
\usepackage[utf8]{inputenc}
\usepackage[scaled=0.96]{helvet}
\renewcommand{\familydefault}{\sfdefault}
\usepackage{array}
\usepackage{enumitem}
\usepackage{fancyhdr}
\usepackage[hidelinks]{hyperref}
\usepackage{microtype}
\usepackage{needspace}
\usepackage{tabularx}
\usepackage{xcolor}
\usepackage{xurl}

\definecolor{aliusgreen}{HTML}{13823A}
\definecolor{aliusgray}{HTML}{7A7A7A}
\definecolor{aliusline}{HTML}{9A9A9A}

\hypersetup{
  colorlinks=true,
  linkcolor=aliusgreen,
  urlcolor=aliusgreen
}

\setlength{\parindent}{0pt}
\setlength{\parskip}{0.72em}
\setlength{\tabcolsep}{0pt}
\raggedbottom
\sloppy
\urlstyle{same}

\newcommand{\ALIUSFooterTitle}{ALIUS Bulletin Guidelines (2022)}

\pagestyle{fancy}
\fancyhf{}
\fancyfoot[L]{\footnotesize\color{aliusgray}\ALIUSFooterTitle}
\fancyfoot[R]{\footnotesize\color{aliusgray}\thepage}
\renewcommand{\headrulewidth}{0pt}
\renewcommand{\footrulewidth}{0.45pt}

\setlist[itemize]{
  leftmargin=1.35em,
  itemsep=0.18em,
  topsep=0.15em,
  parsep=0em
}
\setlist[enumerate]{
  leftmargin=2.05em,
  itemsep=0.42em,
  topsep=0.2em,
  parsep=0em
}
\setlist[description]{
  leftmargin=0em,
  labelsep=0.55em,
  itemsep=0.35em,
  topsep=0.2em,
  font=\normalfont\bfseries
}

\newcommand{\guideDivider}{%
  \par\noindent{\color{aliusline}\rule{\textwidth}{0.45pt}}\par
}

\newcommand{\guideMeta}[2]{%
  {\footnotesize\bfseries\MakeUppercase{#1}}\par
  {\footnotesize #2}\par\vspace{0.65em}
}

\newcommand{\guideHeader}[5]{%
  \thispagestyle{fancy}%
  \vspace*{1mm}
  \begin{center}
    {\Large\bfseries\MakeUppercase{ALIUS Bulletin}\par}
    \vspace{0.25em}
    {\normalsize *\par}
    \vspace{0.25em}
    {\normalsize\bfseries\MakeUppercase{Guidelines For}\par}
    \vspace{0.15em}
    {\LARGE\bfseries\MakeUppercase{#1}\par}
  \end{center}
  \vspace{1.05em}
  \begin{tabularx}{\textwidth}{@{}>{\raggedright\arraybackslash}p{0.55\textwidth}!{\hspace{1.35em}{\color{aliusline}\vrule width 0.5pt}\hspace{1.35em}}>{\raggedright\arraybackslash}X@{}}
    {\large\itshape #2} &
    \guideMeta{Format}{#3}
    \guideMeta{Editorial Route}{#4}
    \guideMeta{Contact}{#5}
  \end{tabularx}
  \vspace{0.75em}
  \guideDivider
  \vspace{0.25em}
}

\newcommand{\guideSection}[1]{%
  \Needspace{6\baselineskip}%
  \par\vspace{1.1em}
  \begin{center}
    {\color{aliusline}\rule{0.14\textwidth}{0.45pt}}\hspace{0.85em}%
    {\large\bfseries #1}%
    \hspace{0.85em}{\color{aliusline}\rule{0.14\textwidth}{0.45pt}}
  \end{center}
  \vspace{-0.25em}
}

\newcommand{\guideSubsection}[1]{%
  \par\vspace{0.65em}{\bfseries #1}\par\vspace{-0.35em}
}

\newcommand{\guideQuestion}[1]{%
  \par\vspace{0.55em}{\color{aliusgreen}\bfseries #1}\par
}

\newenvironment{guideChecklist}
  {\begin{itemize}[label=\textendash]}
  {\end{itemize}}
}
\renewcommand{\ALIUSFooterTitle}{ALIUS Bulletin Guidelines for Podcast (2022)}

\begin{document}

\guideHeader
  {Podcast}
  {If you want to conduct a ``What Is It Like?'' podcast, please read the following guidelines carefully.}
  {Five- to ten-minute episode based on a 500--1,000 word script with bibliographic sources.}
  {Text is edited first; audio, soundtrack, and thumbnail are produced after validation.}
  {\href{mailto:aliusresearchgroup@gmail.com}{aliusresearchgroup@gmail.com}}

\guideSection{Form of the Podcast Episode}

The podcast aims to last between five and ten minutes. The author provides a written text, between 500 and 1,000 words, with bibliographic sources. Editors then turn it into an audio format. It is read by an English speaker with an audio soundtrack added in the background. Timestamps are used for quoting the scientific literature.

\guideSection{Content of the Podcast Episode}

``What is it like to be\ldots?'' intends to provide an overview of the scientific uptake about what it is like to be in a certain conscious state. It favors an explanation of the phenomenological features of this conscious state over mechanistic explanations or functional accounts. It needs to be anchored in the scientific literature, but remains accessible to a large audience and avoids unexplained technical terms.

Personal and anecdotal takes can be provided as long as they cannot be confused with scientific data. Historical insights about how the scientific understanding of the phenomenology of this state has evolved over time are welcome, as is discussion of current research questions and experimental challenges.

A full bibliography and quotations in the text are provided by authors. Please consult the ``what is it like to be asleep'' text and episode as an example if needed.

The podcast is shared on YouTube and podcast platforms. A soundtrack is used to give a subjective rendition of what it is like to be in the given conscious state. The soundtrack will be used to give rhythm to the episode and will be tuned down at certain points to highlight important excerpts of the text.

\guideSection{Modus Operandi and Requirements}

\begin{enumerate}
\item Inform the podcast editors that you want to write about what it is like to be in a conscious state.
\item Send the editors your text, 500--1,000 words, with quotations and full bibliography.
\item Editors proofread your text and suggest edits to improve understanding and flow.
\item Provide a revised edition of the text based on editors' suggestions.
\item Once the text is validated, a speaker is agreed upon by writers and editors and the audio version of the podcast is recorded. No further changes to the text will be possible at this point. The podcast is recorded, the soundtrack is created to fit the audio, and a thumbnail is added.
\item The final version of the podcast is sent to authors for approval.
\item The podcast episode is published depending on the editorial agenda.
\end{enumerate}

\guideSection{Editorial Team}

For any questions, feel free to ask the editors.

L\'ena Coutrot, George Fejer, Daniel Friedman, Matthieu Koroma.

\end{document}
